\chapter*{PHỤ LỤC}
\addcontentsline{toc}{chapter}{PHỤ LỤC}
\markboth{PHỤ LỤC}{}

Phụ lục liệt kê toàn bộ tài nguyên công khai phục vụ việc tái lập (reproducibility) kết quả đề
tài. Tất cả các URL dưới đây có thể truy cập công khai tại thời điểm nộp đề tài.

\section*{Phụ lục A: Mã nguồn (GitHub)}
\addcontentsline{toc}{section}{Phụ lục A: Mã nguồn (GitHub)}
\label{app:github}

\begin{itemize}[leftmargin=1.5cm]
\item \textbf{Thư mục dự án:} \texttt{data-analysis/traffic-yolo-analysis/}
\item \textbf{Cấu trúc thư mục quan trọng:}
\begin{itemize}
\item \texttt{data/notebook/} --- notebook Kaggle nguồn (\texttt{traffic-yolo-v2-run.ipynb})
\item \texttt{src/} --- mã trích xuất, phân tích kiểm định, vẽ hình
\item \texttt{results/tables/} --- bảng CSV trích từ output notebook
\item \texttt{results/figs/} --- hình sinh tự động
\item \texttt{app/app.py} --- ứng dụng web Streamlit
\item \texttt{tests/} --- kiểm thử tự động (pytest)
\item \texttt{reports/} --- báo cáo Word, LaTeX và slide
\end{itemize}
\end{itemize}

\section*{Phụ lục B: Bộ dữ liệu}
\addcontentsline{toc}{section}{Phụ lục B: Bộ dữ liệu}
\label{app:dataset}

\begin{itemize}[leftmargin=1.5cm]
\item \textbf{Bộ dữ liệu:} \url{https://www.kaggle.com/datasets/pkdarabi/cardetection}
\begin{itemize}
\item Quy mô: 4.969 ảnh, 6.012 đối tượng, 15 lớp (đèn tín hiệu + biển báo).
\item Notebook nguồn: \url{https://www.kaggle.com/code/tudeptraine/traffic-yolo-v2-run}
\item Ngày truy cập: tháng 7/2026.
\end{itemize}
\end{itemize}

\section*{Phụ lục C: Hướng dẫn tái lập}
\addcontentsline{toc}{section}{Phụ lục C: Hướng dẫn tái lập}
\label{app:repro}

\begin{lstlisting}[language=bash,caption={Các bước tái lập kết quả}]
# 1. Cai dat moi truong
python -m venv .venv && source .venv/bin/activate
pip install -r requirements.txt

# 2. Chay toan bo pipeline phan tich thu cap
#    (trich xuat -> kiem dinh -> ve hinh)
python -m src.run_pipeline

# 3. Kiem thu tu dong
pytest

# 4. Mo ung dung web
streamlit run app/app.py
\end{lstlisting}

\textbf{Môi trường:} Python 3, các thư viện \texttt{pandas}, \texttt{scipy}, \texttt{matplotlib};
hạt giống ngẫu nhiên \texttt{seed=42}. Số liệu gốc do notebook Kaggle sinh trên GPU Tesla T4
(Ultralytics 8.4.95). Phân tích thứ cấp chạy hoàn toàn trên CPU.
